%% Generated by Sphinx.
\def\sphinxdocclass{report}
\documentclass[a4paper,10pt,oneside,english]{sphinxmanual}
\ifdefined\pdfpxdimen
   \let\sphinxpxdimen\pdfpxdimen\else\newdimen\sphinxpxdimen
\fi \sphinxpxdimen=.75bp\relax
\ifdefined\pdfimageresolution
    \pdfimageresolution= \numexpr \dimexpr1in\relax/\sphinxpxdimen\relax
\fi
%% let collapsible pdf bookmarks panel have high depth per default
\PassOptionsToPackage{bookmarksdepth=5}{hyperref}

\PassOptionsToPackage{booktabs}{sphinx}
\PassOptionsToPackage{colorrows}{sphinx}

\PassOptionsToPackage{warn}{textcomp}
\usepackage[utf8]{inputenc}
\ifdefined\DeclareUnicodeCharacter
% support both utf8 and utf8x syntaxes
  \ifdefined\DeclareUnicodeCharacterAsOptional
    \def\sphinxDUC#1{\DeclareUnicodeCharacter{"#1}}
  \else
    \let\sphinxDUC\DeclareUnicodeCharacter
  \fi
  \sphinxDUC{00A0}{\nobreakspace}
  \sphinxDUC{2500}{\sphinxunichar{2500}}
  \sphinxDUC{2502}{\sphinxunichar{2502}}
  \sphinxDUC{2514}{\sphinxunichar{2514}}
  \sphinxDUC{251C}{\sphinxunichar{251C}}
  \sphinxDUC{2572}{\textbackslash}
\fi
\usepackage{cmap}
\usepackage[T1]{fontenc}
\usepackage{amsmath,amssymb,amstext}
\usepackage[english]{babel}



\usepackage{tgtermes}
\usepackage{tgheros}
\renewcommand{\ttdefault}{txtt}



\usepackage[Bjarne]{fncychap}
\usepackage{sphinx}

\fvset{fontsize=auto}
\usepackage{geometry}


% Include hyperref last.
\usepackage{hyperref}
% Fix anchor placement for figures with captions.
\usepackage{hypcap}% it must be loaded after hyperref.
% Set up styles of URL: it should be placed after hyperref.
\urlstyle{same}

\addto\captionsenglish{\renewcommand{\contentsname}{Contents:}}

\usepackage{sphinxmessages}
\setcounter{tocdepth}{1}



\title{ACSEFunctions Documentation}
\date{Oct 17, 2024}
\release{0.1.0}
\author{Hazem Bakhshwain}
\newcommand{\sphinxlogo}{\vbox{}}
\renewcommand{\releasename}{Release}
\makeindex
\begin{document}

\ifdefined\shorthandoff
  \ifnum\catcode`\=\string=\active\shorthandoff{=}\fi
  \ifnum\catcode`\"=\active\shorthandoff{"}\fi
\fi

\pagestyle{empty}
\sphinxmaketitle
\pagestyle{plain}
\sphinxtableofcontents
\pagestyle{normal}
\phantomsection\label{\detokenize{index::doc}}


\sphinxAtStartPar
Welcome to the \sphinxstylestrong{ACSEFunctions} documentation.

\sphinxstepscope


\chapter{acsefunctions package}
\label{\detokenize{acsefunctions:acsefunctions-package}}\label{\detokenize{acsefunctions::doc}}

\section{Submodules}
\label{\detokenize{acsefunctions:submodules}}

\section{acsefunctions.functions module}
\label{\detokenize{acsefunctions:module-acsefunctions.functions}}\label{\detokenize{acsefunctions:acsefunctions-functions-module}}\index{module@\spxentry{module}!acsefunctions.functions@\spxentry{acsefunctions.functions}}\index{acsefunctions.functions@\spxentry{acsefunctions.functions}!module@\spxentry{module}}
\sphinxAtStartPar
This module contains mathematical functions.

\sphinxAtStartPar
functions include such as factorial, exp, sinh, cosh, and tanh.
\index{bessel() (in module acsefunctions.functions)@\spxentry{bessel()}\spxextra{in module acsefunctions.functions}}

\begin{fulllineitems}
\phantomsection\label{\detokenize{acsefunctions:acsefunctions.functions.bessel}}
\pysigstartsignatures
\pysiglinewithargsret
{\sphinxcode{\sphinxupquote{acsefunctions.functions.}}\sphinxbfcode{\sphinxupquote{bessel}}}
{\sphinxparam{\DUrole{n}{alpha}}\sphinxparamcomma \sphinxparam{\DUrole{n}{x}}\sphinxparamcomma \sphinxparam{\DUrole{n}{N}\DUrole{o}{=}\DUrole{default_value}{100}}}
{}
\pysigstopsignatures
\sphinxAtStartPar
Compute the Bessel function J\_alpha(x) using the series expansion.

\sphinxAtStartPar
This function approximates the Bessel function using a truncated
series expansion
up to N terms.
\begin{quote}\begin{description}
\sphinxlineitem{Parameters}\begin{itemize}
\item {} 
\sphinxAtStartPar
\sphinxstyleliteralstrong{\sphinxupquote{alpha}} (\sphinxstyleliteralemphasis{\sphinxupquote{float}}) \textendash{} The order of the Bessel function (can be any real number).

\item {} 
\sphinxAtStartPar
\sphinxstyleliteralstrong{\sphinxupquote{x}} (\sphinxstyleliteralemphasis{\sphinxupquote{float}}\sphinxstyleliteralemphasis{\sphinxupquote{ or }}\sphinxstyleliteralemphasis{\sphinxupquote{array\sphinxhyphen{}like}}) \textendash{} The value(s) for which to compute the Bessel function.

\item {} 
\sphinxAtStartPar
\sphinxstyleliteralstrong{\sphinxupquote{N}} (\sphinxstyleliteralemphasis{\sphinxupquote{int}}\sphinxstyleliteralemphasis{\sphinxupquote{, }}\sphinxstyleliteralemphasis{\sphinxupquote{optional}}) \textendash{} The number of terms in the series to use for the approximation
(default is 100).

\end{itemize}

\sphinxlineitem{Returns}
\sphinxAtStartPar
The computed value(s) of J\_alpha(x).

\sphinxlineitem{Return type}
\sphinxAtStartPar
np.ndarray

\end{description}\end{quote}
\subsubsection*{Notes}

\sphinxAtStartPar
The Bessel function of the first kind is defined as:
J\_alpha(x) = Σ (from m=0 to \(\infty\)) {[}(\sphinxhyphen{}1)\textasciicircum{}m / (m! * Γ(m + α + 1)){]}
* (x/2)\textasciicircum{}(2m + α)
This implementation uses a truncated series expansionto approximate the
result.

\end{fulllineitems}

\index{cosh() (in module acsefunctions.functions)@\spxentry{cosh()}\spxextra{in module acsefunctions.functions}}

\begin{fulllineitems}
\phantomsection\label{\detokenize{acsefunctions:acsefunctions.functions.cosh}}
\pysigstartsignatures
\pysiglinewithargsret
{\sphinxcode{\sphinxupquote{acsefunctions.functions.}}\sphinxbfcode{\sphinxupquote{cosh}}}
{\sphinxparam{\DUrole{n}{x}}\sphinxparamcomma \sphinxparam{\DUrole{n}{N}\DUrole{o}{=}\DUrole{default_value}{10}}}
{}
\pysigstopsignatures
\sphinxAtStartPar
Compute the hyperbolic cosine of x using the truncated Taylor series.

\sphinxAtStartPar
This function approximates the value of cosh(x) using a
truncated Taylor series
up to N terms. It supports scalar and array inputs, computing the result
element\sphinxhyphen{}wise for arrays.
\begin{quote}\begin{description}
\sphinxlineitem{Parameters}\begin{itemize}
\item {} 
\sphinxAtStartPar
\sphinxstyleliteralstrong{\sphinxupquote{x}} (\sphinxstyleliteralemphasis{\sphinxupquote{float}}\sphinxstyleliteralemphasis{\sphinxupquote{ or }}\sphinxstyleliteralemphasis{\sphinxupquote{array\sphinxhyphen{}like}}) \textendash{} The input value(s) for which to compute cosh(x). Can be a scalar or a
NumPy array.

\item {} 
\sphinxAtStartPar
\sphinxstyleliteralstrong{\sphinxupquote{N}} (\sphinxstyleliteralemphasis{\sphinxupquote{int}}\sphinxstyleliteralemphasis{\sphinxupquote{, }}\sphinxstyleliteralemphasis{\sphinxupquote{optional}}) \textendash{} The number of terms in the Taylor series to use for the approximation
(default is 100). Higher values increase accuracy but also
increase computation time.

\end{itemize}

\sphinxlineitem{Returns}
\sphinxAtStartPar
The computed value(s) of cosh(x) as a NumPy array.
If the input is scalar,
a scalar value is returned; otherwise, an array of the same shape as
the input.

\sphinxlineitem{Return type}
\sphinxAtStartPar
np.ndarray

\end{description}\end{quote}
\subsubsection*{Examples}

\begin{sphinxVerbatim}[commandchars=\\\{\}]
\PYG{g+gp}{\PYGZgt{}\PYGZgt{}\PYGZgt{} }\PYG{n}{cosh}\PYG{p}{(}\PYG{l+m+mi}{0}\PYG{p}{)}
\PYG{g+go}{1.0}
\end{sphinxVerbatim}

\begin{sphinxVerbatim}[commandchars=\\\{\}]
\PYG{g+gp}{\PYGZgt{}\PYGZgt{}\PYGZgt{} }\PYG{n}{cosh}\PYG{p}{(}\PYG{l+m+mi}{1}\PYG{p}{)}
\PYG{g+go}{1.5430806348152437}
\end{sphinxVerbatim}

\begin{sphinxVerbatim}[commandchars=\\\{\}]
\PYG{g+gp}{\PYGZgt{}\PYGZgt{}\PYGZgt{} }\PYG{n}{cosh}\PYG{p}{(}\PYG{n}{np}\PYG{o}{.}\PYG{n}{array}\PYG{p}{(}\PYG{p}{[}\PYG{l+m+mi}{1}\PYG{p}{,} \PYG{l+m+mi}{2}\PYG{p}{,} \PYG{l+m+mi}{3}\PYG{p}{]}\PYG{p}{)}\PYG{p}{)}
\PYG{g+go}{array([ 1.54308063,  3.76219569, 10.06766199])}
\end{sphinxVerbatim}
\subsubsection*{Notes}
\begin{itemize}
\item {} 
\sphinxAtStartPar
The calculation is performed using NumPy’s vectorized

\end{itemize}
\begin{description}
\sphinxlineitem{operations to support}
\sphinxAtStartPar
both scalar and array inputs efficiently.

\end{description}
\begin{itemize}
\item {} 
\sphinxAtStartPar
The function incrementally computes each term in the series using the

\end{itemize}
\begin{description}
\sphinxlineitem{previous}
\sphinxAtStartPar
term, optimizing performance and reducing the risk of overflow.

\end{description}

\end{fulllineitems}

\index{exp() (in module acsefunctions.functions)@\spxentry{exp()}\spxextra{in module acsefunctions.functions}}

\begin{fulllineitems}
\phantomsection\label{\detokenize{acsefunctions:acsefunctions.functions.exp}}
\pysigstartsignatures
\pysiglinewithargsret
{\sphinxcode{\sphinxupquote{acsefunctions.functions.}}\sphinxbfcode{\sphinxupquote{exp}}}
{\sphinxparam{\DUrole{n}{x}}\sphinxparamcomma \sphinxparam{\DUrole{n}{N}\DUrole{o}{=}\DUrole{default_value}{100}}}
{}
\pysigstopsignatures
\sphinxAtStartPar
Compute the exponential of x using the truncated Taylor series.

\sphinxAtStartPar
This function approximates the value of e\textasciicircum{}x using a truncated Taylor series
up to N terms. It supports scalar and array inputs, computing the result
element\sphinxhyphen{}wise for arrays.
\begin{quote}\begin{description}
\sphinxlineitem{Parameters}\begin{itemize}
\item {} 
\sphinxAtStartPar
\sphinxstyleliteralstrong{\sphinxupquote{x}} (\sphinxstyleliteralemphasis{\sphinxupquote{float}}\sphinxstyleliteralemphasis{\sphinxupquote{ or }}\sphinxstyleliteralemphasis{\sphinxupquote{array\sphinxhyphen{}like}}) \textendash{} The input value(s) for which to compute e\textasciicircum{}x.
Can be a scalar or a NumPy array.

\item {} 
\sphinxAtStartPar
\sphinxstyleliteralstrong{\sphinxupquote{N}} (\sphinxstyleliteralemphasis{\sphinxupquote{int}}\sphinxstyleliteralemphasis{\sphinxupquote{, }}\sphinxstyleliteralemphasis{\sphinxupquote{optional}}) \textendash{} The number of terms in the Taylor series to use for the approximation
(default is 100).
Higher values increase accuracy but also increase computation time.

\end{itemize}

\sphinxlineitem{Returns}
\sphinxAtStartPar
The computed value(s) of e\textasciicircum{}x as a NumPy array.
If the input is scalar, a scalar value is returned;
otherwise, an array of the same shape as the input.

\sphinxlineitem{Return type}
\sphinxAtStartPar
np.ndarray

\end{description}\end{quote}
\subsubsection*{Examples}

\begin{sphinxVerbatim}[commandchars=\\\{\}]
\PYG{g+gp}{\PYGZgt{}\PYGZgt{}\PYGZgt{} }\PYG{n}{exp}\PYG{p}{(}\PYG{l+m+mi}{1}\PYG{p}{)}
\PYG{g+go}{2.718281828459045}
\end{sphinxVerbatim}

\begin{sphinxVerbatim}[commandchars=\\\{\}]
\PYG{g+gp}{\PYGZgt{}\PYGZgt{}\PYGZgt{} }\PYG{n}{exp}\PYG{p}{(}\PYG{n}{np}\PYG{o}{.}\PYG{n}{array}\PYG{p}{(}\PYG{p}{[}\PYG{l+m+mi}{1}\PYG{p}{,} \PYG{l+m+mi}{2}\PYG{p}{,} \PYG{l+m+mi}{3}\PYG{p}{]}\PYG{p}{)}\PYG{p}{)}
\PYG{g+go}{array([ 2.71828183,  7.3890561 , 20.08553692])}
\end{sphinxVerbatim}
\subsubsection*{Notes}
\begin{itemize}
\item {} 
\sphinxAtStartPar
The calculation is performed using NumPy’s

\end{itemize}
\begin{quote}

\sphinxAtStartPar
vectorized operations to support both scalar and array inputs efficiently.
\end{quote}
\begin{itemize}
\item {} 
\sphinxAtStartPar
The function incrementally computes each term in the series using the
previous term,
which helps minimize the risk of overflow and maximizes efficiency.

\end{itemize}

\end{fulllineitems}

\index{fact() (in module acsefunctions.functions)@\spxentry{fact()}\spxextra{in module acsefunctions.functions}}

\begin{fulllineitems}
\phantomsection\label{\detokenize{acsefunctions:acsefunctions.functions.fact}}
\pysigstartsignatures
\pysiglinewithargsret
{\sphinxcode{\sphinxupquote{acsefunctions.functions.}}\sphinxbfcode{\sphinxupquote{fact}}}
{\sphinxparam{\DUrole{n}{x}}}
{}
\pysigstopsignatures
\sphinxAtStartPar
Calculate the factorial of a non\sphinxhyphen{}negative integer or NumPyarray.
\begin{quote}\begin{description}
\sphinxlineitem{Parameters}
\sphinxAtStartPar
\sphinxstyleliteralstrong{\sphinxupquote{x}} (\sphinxstyleliteralemphasis{\sphinxupquote{int}}\sphinxstyleliteralemphasis{\sphinxupquote{ or }}\sphinxstyleliteralemphasis{\sphinxupquote{np.ndarray}}) \textendash{} The non\sphinxhyphen{}negative integer or array of integers for which the factorial
is computed.

\sphinxlineitem{Returns}
\sphinxAtStartPar
The factorial of x (x!) if x is a scalar, or an array of factorials
if x is a NumPy array.

\sphinxlineitem{Return type}
\sphinxAtStartPar
int or np.ndarray

\end{description}\end{quote}
\subsubsection*{Examples}

\begin{sphinxVerbatim}[commandchars=\\\{\}]
\PYG{g+gp}{\PYGZgt{}\PYGZgt{}\PYGZgt{} }\PYG{n}{fact}\PYG{p}{(}\PYG{l+m+mi}{5}\PYG{p}{)}
\PYG{g+go}{120}
\end{sphinxVerbatim}

\begin{sphinxVerbatim}[commandchars=\\\{\}]
\PYG{g+gp}{\PYGZgt{}\PYGZgt{}\PYGZgt{} }\PYG{n}{fact}\PYG{p}{(}\PYG{n}{np}\PYG{o}{.}\PYG{n}{array}\PYG{p}{(}\PYG{p}{[}\PYG{l+m+mi}{0}\PYG{p}{,} \PYG{l+m+mi}{1}\PYG{p}{,} \PYG{l+m+mi}{2}\PYG{p}{,} \PYG{l+m+mi}{3}\PYG{p}{,} \PYG{l+m+mi}{4}\PYG{p}{]}\PYG{p}{)}\PYG{p}{)}
\PYG{g+go}{array([ 1,  1,  2,  6, 24])}
\end{sphinxVerbatim}
\subsubsection*{Notes}

\sphinxAtStartPar
Uses recursion and caching (@cache) for efficiency in scalar calculations.
For NumPy arrays, the function is vectorized to handle each element.

\end{fulllineitems}

\index{gamma() (in module acsefunctions.functions)@\spxentry{gamma()}\spxextra{in module acsefunctions.functions}}

\begin{fulllineitems}
\phantomsection\label{\detokenize{acsefunctions:acsefunctions.functions.gamma}}
\pysigstartsignatures
\pysiglinewithargsret
{\sphinxcode{\sphinxupquote{acsefunctions.functions.}}\sphinxbfcode{\sphinxupquote{gamma}}}
{\sphinxparam{\DUrole{n}{z}}}
{}
\pysigstopsignatures
\sphinxAtStartPar
Compute the gamma function Γ(z) using the Lanczos approximation.

\sphinxAtStartPar
This function uses the Lanczos approximation, which is known for its
accuracy
and efficiency in computing the gamma function.
\begin{quote}\begin{description}
\sphinxlineitem{Parameters}
\sphinxAtStartPar
\sphinxstyleliteralstrong{\sphinxupquote{z}} (\sphinxstyleliteralemphasis{\sphinxupquote{float}}) \textendash{} The input value for which to compute the gamma function.
Should be positive.

\sphinxlineitem{Returns}
\sphinxAtStartPar
The computed value of Γ(z).

\sphinxlineitem{Return type}
\sphinxAtStartPar
float

\end{description}\end{quote}
\subsubsection*{Notes}
\begin{description}
\sphinxlineitem{The Lanczos approximation formula is:}
\sphinxAtStartPar
Γ(z) ≈ sqrt(2π) * (z + g \sphinxhyphen{} 0.5)\textasciicircum{}(z \sphinxhyphen{} 0.5) * exp(\sphinxhyphen{}(z + g \sphinxhyphen{} 0.5)) * A(z)

\end{description}

\sphinxAtStartPar
where A(z) is a weighted sum of coefficients.

\end{fulllineitems}

\index{sinh() (in module acsefunctions.functions)@\spxentry{sinh()}\spxextra{in module acsefunctions.functions}}

\begin{fulllineitems}
\phantomsection\label{\detokenize{acsefunctions:acsefunctions.functions.sinh}}
\pysigstartsignatures
\pysiglinewithargsret
{\sphinxcode{\sphinxupquote{acsefunctions.functions.}}\sphinxbfcode{\sphinxupquote{sinh}}}
{\sphinxparam{\DUrole{n}{x}}\sphinxparamcomma \sphinxparam{\DUrole{n}{N}\DUrole{o}{=}\DUrole{default_value}{10}}}
{}
\pysigstopsignatures
\sphinxAtStartPar
Compute the hyperbolic sine of x using the truncated Taylor series.

\sphinxAtStartPar
This function approximates the value of sinh(x)
using a truncated Taylor series
up to N terms. It supports scalar and array inputs, computing the result
element\sphinxhyphen{}wise for arrays.
\begin{quote}\begin{description}
\sphinxlineitem{Parameters}\begin{itemize}
\item {} 
\sphinxAtStartPar
\sphinxstyleliteralstrong{\sphinxupquote{x}} (\sphinxstyleliteralemphasis{\sphinxupquote{float}}\sphinxstyleliteralemphasis{\sphinxupquote{ or }}\sphinxstyleliteralemphasis{\sphinxupquote{array\sphinxhyphen{}like}}) \textendash{} The input value(s) for which to compute sinh(x). Can be a scalar or a
NumPy array.

\item {} 
\sphinxAtStartPar
\sphinxstyleliteralstrong{\sphinxupquote{N}} (\sphinxstyleliteralemphasis{\sphinxupquote{int}}\sphinxstyleliteralemphasis{\sphinxupquote{, }}\sphinxstyleliteralemphasis{\sphinxupquote{optional}}) \textendash{} The number of terms in the Taylor series to use for the approximation
(default is 100). Higher values increase accuracy but also increase
computation time.

\end{itemize}

\sphinxlineitem{Returns}
\sphinxAtStartPar
The computed value(s) of sinh(x) as a NumPy array.
If the input is scalar,
a scalar value is returned; otherwise,
an array of the same shape as the input.

\sphinxlineitem{Return type}
\sphinxAtStartPar
np.ndarray

\end{description}\end{quote}
\subsubsection*{Examples}

\begin{sphinxVerbatim}[commandchars=\\\{\}]
\PYG{g+gp}{\PYGZgt{}\PYGZgt{}\PYGZgt{} }\PYG{n}{sinh}\PYG{p}{(}\PYG{l+m+mi}{0}\PYG{p}{)}
\PYG{g+go}{0.0}
\end{sphinxVerbatim}

\begin{sphinxVerbatim}[commandchars=\\\{\}]
\PYG{g+gp}{\PYGZgt{}\PYGZgt{}\PYGZgt{} }\PYG{n}{sinh}\PYG{p}{(}\PYG{l+m+mi}{1}\PYG{p}{)}
\PYG{g+go}{1.1752011936438014}
\end{sphinxVerbatim}

\begin{sphinxVerbatim}[commandchars=\\\{\}]
\PYG{g+gp}{\PYGZgt{}\PYGZgt{}\PYGZgt{} }\PYG{n}{sinh}\PYG{p}{(}\PYG{n}{np}\PYG{o}{.}\PYG{n}{array}\PYG{p}{(}\PYG{p}{[}\PYG{l+m+mi}{1}\PYG{p}{,} \PYG{l+m+mi}{2}\PYG{p}{,} \PYG{l+m+mi}{3}\PYG{p}{]}\PYG{p}{)}\PYG{p}{)}
\PYG{g+go}{array([ 1.17520119,  3.62686041, 10.01787493])}
\end{sphinxVerbatim}
\subsubsection*{Notes}
\begin{itemize}
\item {} 
\sphinxAtStartPar
The calculation is performed using NumPy’s vectorized operations to

\end{itemize}
\begin{description}
\sphinxlineitem{support}
\sphinxAtStartPar
both scalar and array inputs efficiently.

\end{description}
\begin{itemize}
\item {} 
\sphinxAtStartPar
The function incrementally computes each term in the series using the

\end{itemize}
\begin{description}
\sphinxlineitem{previous}
\sphinxAtStartPar
power of x, which helps minimize computational overhead.

\end{description}

\end{fulllineitems}

\index{tanh() (in module acsefunctions.functions)@\spxentry{tanh()}\spxextra{in module acsefunctions.functions}}

\begin{fulllineitems}
\phantomsection\label{\detokenize{acsefunctions:acsefunctions.functions.tanh}}
\pysigstartsignatures
\pysiglinewithargsret
{\sphinxcode{\sphinxupquote{acsefunctions.functions.}}\sphinxbfcode{\sphinxupquote{tanh}}}
{\sphinxparam{\DUrole{n}{x}}\sphinxparamcomma \sphinxparam{\DUrole{n}{N}\DUrole{o}{=}\DUrole{default_value}{10}}}
{}
\pysigstopsignatures
\sphinxAtStartPar
Compute the hyperbolic tangent of x using sinh(x) and cosh(x).

\sphinxAtStartPar
This function calculates tanh(x) as the ratio of sinh(x) and cosh(x), using
a truncated Taylor series up to N terms.
It supports scalar and array inputs,
computing the result element\sphinxhyphen{}wise for arrays.
\begin{quote}\begin{description}
\sphinxlineitem{Parameters}\begin{itemize}
\item {} 
\sphinxAtStartPar
\sphinxstyleliteralstrong{\sphinxupquote{x}} (\sphinxstyleliteralemphasis{\sphinxupquote{float}}\sphinxstyleliteralemphasis{\sphinxupquote{ or }}\sphinxstyleliteralemphasis{\sphinxupquote{array\sphinxhyphen{}like}}) \textendash{} The input value(s) for which to compute tanh(x). Can be a scalar or a
NumPy array.

\item {} 
\sphinxAtStartPar
\sphinxstyleliteralstrong{\sphinxupquote{N}} (\sphinxstyleliteralemphasis{\sphinxupquote{int}}\sphinxstyleliteralemphasis{\sphinxupquote{, }}\sphinxstyleliteralemphasis{\sphinxupquote{optional}}) \textendash{} The number of terms in the Taylor series used for sinh(x) and cosh(x)
approximations (default is 100). Higher values increase accuracy but
also increase computation time.

\end{itemize}

\sphinxlineitem{Returns}
\sphinxAtStartPar
The computed value(s) of tanh(x) as a NumPy array. If the input
is scalar,
a scalar value is returned; otherwise, an array of the same shape as
the input.

\sphinxlineitem{Return type}
\sphinxAtStartPar
np.ndarray

\end{description}\end{quote}
\subsubsection*{Examples}

\begin{sphinxVerbatim}[commandchars=\\\{\}]
\PYG{g+gp}{\PYGZgt{}\PYGZgt{}\PYGZgt{} }\PYG{n}{tanh}\PYG{p}{(}\PYG{l+m+mi}{0}\PYG{p}{)}
\PYG{g+go}{0.0}
\end{sphinxVerbatim}

\begin{sphinxVerbatim}[commandchars=\\\{\}]
\PYG{g+gp}{\PYGZgt{}\PYGZgt{}\PYGZgt{} }\PYG{n}{tanh}\PYG{p}{(}\PYG{l+m+mi}{1}\PYG{p}{)}
\PYG{g+go}{0.7615941559557649}
\end{sphinxVerbatim}

\begin{sphinxVerbatim}[commandchars=\\\{\}]
\PYG{g+gp}{\PYGZgt{}\PYGZgt{}\PYGZgt{} }\PYG{n}{tanh}\PYG{p}{(}\PYG{n}{np}\PYG{o}{.}\PYG{n}{array}\PYG{p}{(}\PYG{p}{[}\PYG{l+m+mi}{1}\PYG{p}{,} \PYG{l+m+mi}{2}\PYG{p}{,} \PYG{l+m+mi}{3}\PYG{p}{]}\PYG{p}{)}\PYG{p}{)}
\PYG{g+go}{array([0.76159416, 0.96402758, 0.99505475])}
\end{sphinxVerbatim}
\subsubsection*{Notes}
\begin{itemize}
\item {} 
\sphinxAtStartPar
This function uses the sinh(x) and cosh(x) functions previously defined,
which are based on truncated Taylor series expansions.

\item {} 
\sphinxAtStartPar
The calculation is performed using NumPy’s vectorized operations to

\end{itemize}
\begin{description}
\sphinxlineitem{support}
\sphinxAtStartPar
both scalar and array inputs efficiently.

\end{description}

\end{fulllineitems}



\section{Module contents}
\label{\detokenize{acsefunctions:module-acsefunctions}}\label{\detokenize{acsefunctions:module-contents}}\index{module@\spxentry{module}!acsefunctions@\spxentry{acsefunctions}}\index{acsefunctions@\spxentry{acsefunctions}!module@\spxentry{module}}
\sphinxAtStartPar
The acsefunctions package provides mathematical functions.

\sphinxAtStartPar
These functions include factorial, exponential, and hyperbolic functions.

\sphinxstepscope


\chapter{acsefunctions}
\label{\detokenize{modules:acsefunctions}}\label{\detokenize{modules::doc}}

\chapter{Indices and tables}
\label{\detokenize{index:indices-and-tables}}\begin{itemize}
\item {} 
\sphinxAtStartPar
\DUrole{xref}{\DUrole{std}{\DUrole{std-ref}{genindex}}}

\item {} 
\sphinxAtStartPar
\DUrole{xref}{\DUrole{std}{\DUrole{std-ref}{modindex}}}

\item {} 
\sphinxAtStartPar
\DUrole{xref}{\DUrole{std}{\DUrole{std-ref}{search}}}

\end{itemize}


\renewcommand{\indexname}{Python Module Index}
\begin{sphinxtheindex}
\let\bigletter\sphinxstyleindexlettergroup
\bigletter{a}
\item\relax\sphinxstyleindexentry{acsefunctions}\sphinxstyleindexpageref{acsefunctions:\detokenize{module-acsefunctions}}
\item\relax\sphinxstyleindexentry{acsefunctions.functions}\sphinxstyleindexpageref{acsefunctions:\detokenize{module-acsefunctions.functions}}
\end{sphinxtheindex}

\renewcommand{\indexname}{Index}
\printindex
\end{document}